% !TeX root = booklet-02.tex
\clearpage
\subsection{جمع‌بندی اثرهای غیرایده‌آل و کاربرد در مدار}
\label{sec:nonideal-summary}
بخش \secnum{۲.۴.۷} کتاب، نتیجه بحث را در پنج موضوع مدار جمع می‌کند
\cite[صفحات ۸۶ و ۸۷]{weste}. اسلاید ۳۰ نیز انتخاب تغذیه،
\lr{Logical Effort}، توان سکون، \lr{Pass Transistor} و دمای کار را
از کاربردهای این اثرها می‌شمارد \cite{nonideal}.

\subsubsection*{افت آستانه در عبور سطح منطقی}
وقتی \lr{nMOS} با گیت \(V_{DD}\) یک گره را به سمت سطح بالا شارژ
می‌کند، افزایش ولتاژ گره، \(V_{GS}\) را کاهش می‌دهد. در تقریب
آستانه ثابت، جریان هدایت قوی نزدیک
\(V_{\mathrm{out}}=V_{DD}-V_{tn}\) ضعیف می‌شود. در \lr{pMOS} که
سطح پایین را عبور می‌دهد، حد متناظر حدود \(|V_{tp}|\) بالای زمین
است. \lr{Body Effect} بزرگی افت آستانه را بیشتر می‌کند. وقتی آستانه
سهم بزرگی از تغذیه باشد، این افت برای سطح منطقی اهمیت بیشتری دارد.
در \lr{Transmission Gate}، یک \lr{nMOS} و یک \lr{pMOS} با کنترل
مکمل به‌صورت موازی قرار می‌گیرند. \lr{nMOS} عبور سطح پایین و
\lr{pMOS} عبور سطح بالا را تقویت می‌کند
\cite[بخش \secnum{۲.۴.۷}]{weste}. استخراج مشخصه انتقال و سطوح
منطقی این مدارها در بخش \secnum{۲.۵}، موضوع جزوه سوم است.

\subsubsection*{نشتی، توان سکون و نگهداری داده}
در \lr{Static CMOS} ایده‌آل، در حالت ماندگار مسیر هدایت بین تغذیه
و زمین قطع است. مدل واقعی مسیرهای نشتی بخش \ref{sec:leakage} را
دارد. در صورت تغذیه واحد و با قرارداد جریان مثبت دریافتی از منبع،
توان ماندگار \(P_{\mathrm{static}}=V_{DD}I_{DD}\) است.
برای چند منبع، سهم هر منبع از حاصل‌ضرب ولتاژ و جریان همان منبع
محاسبه می‌شود. جریان‌های شاخه‌ای باید با \lr{KCL} به جریان منبع
تبدیل شوند.

در گره‌ای که داده را به شکل بار ذخیره می‌کند، نشتی بار را جابه‌جا
می‌کند و ولتاژ گره تغییر می‌کند. برای ظرفیت ثابت،
\(\Delta Q=C\Delta V\) است. اگر جریان نشتی را در یک بازه تقریباً
ثابت فرض کنیم، \(|\Delta Q|\approx|I_{\mathrm{leak}}|\Delta t\)
و زمان رسیدن به تغییر ولتاژ مجاز تقریباً برابر می‌شود با
\begin{equation}
 \Delta t\approx\frac{C|\Delta V|}{|I_{\mathrm{leak}}|}.
 \label{eq:retention-time-estimate}
\end{equation}
واحد \(\mathrm{F}\,\mathrm{V}/\mathrm{A}=\mathrm{s}\) است.
این برآورد به ثابت‌بودن تقریبی ظرفیت و جریان تکیه دارد. در حالت
کلی، تغییر ولتاژ با جریان وابسته به ولتاژ و دما محاسبه می‌شود.
کتاب نیاز به \lr{Feedback} یا تجدید داده در ساختارهای ذخیره‌ساز را
از پیامدهای نشتی مطرح می‌کند. آستانه پایین، جریان روشن و سرعت را
تقویت می‌کند و آستانه بالا نشتی را کاهش می‌دهد. فرایندهای دارای
چند انتخاب آستانه از این وابستگی برای مسیرهای با نیازهای متفاوت
استفاده می‌کنند \cite[صفحه ۸۷]{weste}.

در روایت فناوری مرجع، نازک‌شدن عایق، نشتی تونلی گیت را مهم کرده
و استفاده از عایق‌های \lr{High-k} را برانگیخته است. همچنین شیب
زیرآستانه در نمونه‌های بحث کتاب در حدود
\(\qty{60}{\milli\volt}\) تا \(\qty{100}{\milli\volt}\) به ازای هر
دهه تغییر جریان است. رابطه \eqref{eq:leakage-subthreshold-slope}
وابستگی آن را به \(n\) و دما نشان می‌دهد
\cite[صفحات ۸۲ و ۸۷]{weste}.

\subsubsection*{انتخاب \lr{VDD} و بازده کاهش‌یابنده افزایش ولتاژ}
در مدل مربعی، جریان اشباع با مربع اضافه‌ولتاژ رشد می‌کند. با
\lr{Mobility Degradation} و \lr{Velocity Saturation}، رشد جریان
در بایاس بالا کندتر می‌شود. در نتیجه، افزایش تغذیه با بهبود جریان
کمتری نسبت به پیش‌بینی ساده مربعی همراه است. هم‌زمان، انرژی
سوئیچ‌کردن ظرفیت با \(V_{DD}^2\) رشد می‌کند. تنش عایق نازک و
کانال کوتاه نیز محدوده مجاز ولتاژ را محدود می‌کند
\cite[صفحه ۸۷]{weste}. مقدار مجاز ولتاژ و شرایط قابلیت اطمینان از
اطلاعات همان فناوری تعیین می‌شوند.

\subsubsection*{ترانزیستورهای سری، جریان و تأخیر}
در ساختار سری، ولتاژ کل بین چند افزاره تقسیم می‌شود. کاهش میدان
هر افزاره می‌تواند اثر \lr{Velocity Saturation} را کم کند. کتاب
بیان می‌کند که جریان دو \lr{nMOS} سری می‌تواند از نصف جریان یک
\lr{nMOS} با همان عرض و طول بیشتر باشد. این اثر برای \lr{nMOS}
با \lr{Mobility} بالاتر محسوس‌تر است
\cite[صفحه ۸۷]{weste}. در برآورد
\(t\approx C\Delta V/I\)، جریان بیشتر به زمان کوتاه‌تر برای
جابه‌جایی همان بار کمک می‌کند. مدل دقیق تأخیر به تغییر جریان و
ظرفیت طی گذار وابسته است.

تمرین‌های \secnum{۲.۲} و \secnum{۲.۳}، ساختار سری را با یک ترانزیستور طول \(2L\)
مقایسه می‌کنند. مثال آموزشی بخش \ref{sec:series-velocity-example}
مقایسه کتاب در جمع‌بندی را با یک ترانزیستور طول \(L\) انجام می‌دهد.
مشخص‌کردن افزاره مرجع و مدل، برای تفسیر نسبت جریان ضروری است.

\subsubsection*{تطبیق جریان و \lr{Matching}}
برای برابری رفتار، ابعاد، جهت‌گیری و محیط ساخت افزاره‌ها باید تا
حد امکان مشابه باشند. واحدهای موازی یکسان و درهم‌گذاری مناسب
می‌توانند به ساخت نسبت جریان کمک کنند. مثال بخش
\ref{sec:geometry-dependence} نشان داد که یک نسبت در ابعاد
ترسیم‌شده الزاماً همان نسبت را در ابعاد مؤثر ایجاد نمی‌کند
\cite[صفحات ۸۶ و ۸۷]{weste}.

\subsubsection*{مثال آموزشی جریان دو افزاره سری در مدل سرعت}
\label{sec:series-velocity-example}
دو \lr{nMOS} یکسان با گیت مشترک \(V_{DD}=\qty{1}{\volt}\)،
سورس پایین در زمین و درین بالا در \(\qty{1}{\volt}\) را در نظر
بگیرید. برای هر واحد فرض کنید
\(V_t=\qty{0.3}{\volt}\)، \(V_c=\qty{1}{\volt}\) و
\(\beta=\qty{400}{\micro\ampere\per\volt\squared}\) باشند.
اثر بدنه، \lr{DIBL} و \lr{CLM} در این مثال آموزشی حذف شده‌اند.
ولتاژ گره میانی را \(x\) می‌نامیم.
رابطه‌های \eqref{eq:high-field-linear} و
\eqref{eq:current-reduction-factor} را برای هر افزاره جداگانه به
کار می‌بریم. برای اختصار، اگر \(a=V_{GS}-V_t>0\) و \(d=V_{DS}\)
باشند، تابع جریان همین مدل را می‌نویسیم
\begin{equation}
 J(a,d)=
 \begin{cases}
 \displaystyle\beta\frac{ad-d^2/2}{1+d/V_c},
       &0\le d<\dfrac{aV_c}{a+V_c},\\[6pt]
 \displaystyle\frac{\beta a^2}{2(1+a/V_c)},
       &d\ge\dfrac{aV_c}{a+V_c}.
 \end{cases}
 \label{eq:series-velocity-function}
\end{equation}
برای افزاره پایین، \(a=0.7\) و \(d=x\) برحسب ولت‌اند.
برای افزاره بالا، \(a=0.7-x\) و \(d=1-x\) هستند. شرط برابری
جریان‌ها در گره میانی می‌شود
\begin{equation}
 J(0.7,x)=J(0.7-x,1-x).
 \label{eq:series-velocity-kcl}
\end{equation}
تابع‌ها با ولتاژ برحسب ولت و ضریب \(\beta\) داده‌شده ارزیابی
می‌شوند. ابتدا \(0<x<0.7\) را انتخاب می‌کنیم تا هر دو افزاره
روشن باشند. در \(x=\qty{0.1}{\volt}\)، جریان پایین
\(\qty{23.64}{\micro\ampere}\) و جریان بالا
\(\qty{45.00}{\micro\ampere}\) است. پس باید \(x\) افزایش یابد.
در \(x=\qty{0.2}{\volt}\)، این جریان‌ها به‌ترتیب
\(\qty{40.00}{\micro\ampere}\) و
\(\qty{33.33}{\micro\ampere}\) هستند. ریشه میان این دو مقدار است.
با نصف‌کردن مکرر بازه و انتخاب نیمه‌ای که اختلاف جریان در دو سر
آن تغییر علامت می‌دهد، نتیجه چنین به دست می‌آید
\begin{equation}
 x\approx\qty{0.17370}{\volt},\qquad
 I_{\mathrm{series}}\approx\qty{36.296}{\micro\ampere}.
 \label{eq:series-velocity-result}
\end{equation}
کنترل ناحیه کاری ضروری است. برای افزاره پایین، مرز اشباع
\(0.7/(1+0.7)\approx\qty{0.41176}{\volt}\) است و
\(x\) از آن کمتر است، پس شاخه \lr{Linear} انتخاب می‌شود.
برای افزاره بالا، اضافه‌ولتاژ
\(\qty{0.52630}{\volt}\)، ولتاژ درین و سورس
\(\qty{0.82630}{\volt}\) و مرز اشباع
\(\qty{0.34482}{\volt}\) است. پس افزاره بالا در اشباع است.

یک افزاره واحد طول \(L\) در همان بایاس کل، در اشباع قرار دارد و
جریان آن برابر است با
\begin{equation}
 I_{\mathrm{single}}=
 \frac{(400/2)(0.7)^2}{1+0.7/1}\,\unit{\micro\ampere}
 \approx\qty{57.647}{\micro\ampere}.
 \label{eq:series-velocity-reference}
\end{equation}
بنابراین، \(I_{\mathrm{series}}/I_{\mathrm{single}}\approx0.6296\)
است که از نصف بزرگ‌تر است. این مقدار، نتیجه پارامترها و فرض‌های
همین مثال است. برای افزودن اثر بدنه، باید آستانه افزاره بالا را
تابع \(x\) کرد و برابری جریان‌ها را دوباره حل کرد.

\subsubsection*{تمرین آموزشی نگهداری بار و پاسخ}
گرهی با ظرفیت ثابت \(\qty{10}{\femto\farad}\)، جریان نشتی تقریباً
ثابت \(\qty{1}{\nano\ampere}\) و تغییر ولتاژ مجاز
\(\qty{0.1}{\volt}\) فرض کنید. زمان برآوردی نگهداری بار چقدر است؟
با تبدیل واحدها و استفاده از \eqref{eq:retention-time-estimate}، زمان
برآوردی به دست می‌آید.
\[
 \Delta t=\frac{(10\times10^{-15})(0.1)}{10^{-9}}
 =\qty{1e-6}{\second}=\qty{1}{\micro\second}.
\]
اگر جریان در همان شرایط دو برابر شود و ظرفیت و حاشیه ولتاژ ثابت
بمانند، زمان به \(\qty{0.5}{\micro\second}\) کاهش می‌یابد.
برای نمونه واقعی با جریان متغیر، معادله تغییر بار باید با مدل
جریان همان گره حل شود.

\subsubsection*{راهنمای نهایی حل و نقشه پوشش منابع}
پیش از جای‌گذاری عددها، بایاس‌های پایانه‌ای، مدل مسئله و شرایط ثابت
را تعیین کنید. ابعاد مؤثر را با قرارداد مدل به کار ببرید. آستانه
متناسب با بدنه، درین، هندسه و دما را محاسبه کنید. سپس اضافه‌ولتاژ
و مرز ناحیه کاری را به دست آورید. در مدار چندترانزیستوری، ولتاژ
گره‌های داخلی از \lr{KCL} تعیین می‌شود. پس از حل، ناحیه کاری،
ریشه فیزیکی، واحدها، رفتار حدی و دامنه تقریب‌ها را بررسی کنید.
پارامترهای تجربی چند مدل باید با قرارداد یکسان ترکیب شوند.

جدول \ref{tab:booklet-two-coverage} مسیر بازگشت از موضوع به کتاب و
اسلاید را نشان می‌دهد. شماره‌های رابطه این جزوه مستقل از کتاب‌اند.
\begin{table}[!htbp]
\centering
\small
\caption{نقشه موضوعی بخش \secnum{۲.۴} کتاب و اسلایدهای \lr{Nonideal Transistor Theory}.}
\label{tab:booklet-two-coverage}
\begin{tabularx}{\linewidth}{c>{\raggedleft\arraybackslash}Xc}
\toprule
بخش کتاب & موضوع در همین جزوه & اسلاید\\
\midrule
\secnum{۲.۴.۱} & بخش \ref{sec:high-field}، میدان، سرعت و مدل توان & ۶ تا ۱۲\\
\secnum{۲.۴.۲} & بخش \ref{sec:clm}، \lr{CLM} و شیب خروجی & ۱۳ و ۱۴\\
\secnum{۲.۴.۳} & بخش \ref{sec:threshold-effects}، اثرهای آستانه & ۱۵ تا ۲۰\\
\secnum{۲.۴.۴} & بخش \ref{sec:leakage}، سه مسیر نشتی & ۲۱ تا ۲۸\\
\secnum{۲.۴.۵} & بخش \ref{sec:temperature-dependence}، رقابت اثرهای دما & ۲۹\\
\secnum{۲.۴.۶} & بخش \ref{sec:geometry-dependence}، ابعاد مؤثر و تطبیق & ۳۵ و ۳۶\\
\secnum{۲.۴.۷} & بخش \ref{sec:nonideal-summary}، پیامدهای مداری & ۳۰\\
مکمل دوره & \hyperref[sec:pvt-supplement]{مکمل \lr{PVT} و قراردادهای \lr{Corner}} & ۳۱ تا ۳۴\\
\bottomrule
\end{tabularx}
\end{table}

اسلایدهای ۳ تا ۵ مقدمه و مرور منحنی‌ها هستند و موضوع آن‌ها در مرور
پیش‌نیازها و آغاز بخش \ref{sec:nonideal-current} آمده است.
فهرست اسلاید ۳۷ با هفت تمرین سری نخست فصل دوم مطابقت دارد.
اثبات تمرین \secnum{۲.۶} در بخش \ref{sec:high-field}، حل \secnum{۲.۲} در بازسازی
مرجع مدار سری و حل \secnum{۲.۳} و \secnum{۲.۷} تا \secnum{۲.۹} در فصل تمرین‌های اثر آستانه
آمده است. فصل پایانی تمرین‌ها، مسئله \secnum{۲.۱۰} را با مقایسه مدل مربعی
و نشتی زیرآستانه حل می‌کند. این ترتیب، پیش‌نیاز هر تمرین را پیش از
حل آن فراهم می‌کند.
